\documentclass[10pt,a4paper]{amsart}
\usepackage[margin=19mm]{geometry}
\usepackage{amsmath,amssymb,mathtools}
\usepackage[scr=rsfs,cal=euler]{mathalfa}
\usepackage{xfrac}
\usepackage[unicode,hidelinks,bookmarks=false]{hyperref}
\newcommand{\ie}{i.e.\ }
\newcommand{\cf}{cf.\ }
\newcommand{\resp}{respectively\ }
\newcommand{\Subsection}{\subsection}
\providecommand{\nobreakdash}{\nobreak\hbox{-}}
\setlength{\emergencystretch}{2em}
\multlinegap0pt
\usepackage{bm}
\usepackage{bbm}
\usepackage{dsfont}
\usepackage{xspace}
\usepackage{cancel}
\usepackage{mathtools}
\usepackage{amsthm}
\makeatletter
\@ifundefined{thm}{\newtheorem{thm}{Thm}}{}
\@ifundefined{lem}{\newtheorem{lem}[thm]{Lemma}}{}
\makeatother
\newcommand{\R}{{\mathbb R}}
\newcommand{\Tr}{\operatorname{Tr}}
\newcommand{\cF}{{\mathcal F}}
\newcommand{\cM}{{\mathcal M}}
\newcommand{\cO}{{\mathcal O}}
\newcommand{\ot}{\otimes}
\newcommand{\sH}{\mathsf{H}}
\title[Structural and non-isomorphism results for $q$-Araki-Woods f]{Structural and non-isomorphism results for $q$-Araki-Woods factors}
\author{Changying Ding, Hui Tan}
\date{Source: arXiv:2509.21636v1 (CC BY 4.0)}
\begin{document}
\maketitle

\noindent\emph{Source and licence.} Structural and non-isomorphism results for $q$-Araki-Woods factors. Version arXiv:2509.21636v1 (CC BY 4.0), distributed under the Creative Commons Attribution 4.0 International licence, as recorded in the arXiv licence metadata for this exact version and shown on the abstract page https://arxiv.org/abs/2509.21636v1. Full rights notice: licenses/arxiv-2509.21636-CC-BY-4.0.txt.\par\smallskip

\noindent\textit{Editorial scope.} Counted unit: the Lem whose source label is \texttt{lem: embedding bimodules} in the article (source0.tex lines 734--750, source label \texttt{lem: embedding bimodules}), delivered together with the article's own complete original proof of it (source0.tex lines 751--789, one \texttt{proof} environment of 39 lines). No mathematics is written, repaired or re-derived: statement and proof are the article's own text line for line. Required context retained uncounted: none. Every definition, display and label of those lines is retained unchanged. Omitted from this package and not redistributed: 1--733, 790--1555. Nothing inside a retained excerpt is omitted or altered: every retained body is delivered exactly as the article's lines are written, so no omission receipt of any kind accompanies this package. Citations: the retained excerpts carry no citation command at all, so no bibliography entry of the article is transcribed and no reference is redistributed. Reference adaptation: none was needed and none was made; every reference inside the delivered unit resolves inside it, so no unresolved reference or citation is produced. Printed slips kept literal: no printed word, symbol, hypothesis or number of the counted statement or of its proof was altered. Layout and class: the article's own class is \texttt{amsart} with packages including amsmath, amsthm, amsfonts, amssymb, enumerate, hyperref, xcolor, verbatim, ulem, tikz-cd; the wrapper is the lane's pinned \texttt{amsart} editorial wrapper at 10pt on A4, which restates the theorem-like environments the retained text needs and adds the article's own macro definitions that the retained text needs (\texttt{\textbackslash R}, \texttt{\textbackslash Tr}, \texttt{\textbackslash cF}, \texttt{\textbackslash cM}, \texttt{\textbackslash cO}, \texttt{\textbackslash ot}, \texttt{\textbackslash sH}), with extra wrapper packages (bm, bbm, dsfont, xspace, cancel, mathtools). The delivered numbering is the unit's own. Assets: no retained span contains a figure environment, a graphics-inclusion command, a PDF-page inclusion, a table or a TikZ picture; no image file is copied into this package, the package declares no asset dependency and no image file is redistributed with it. Licence: version arXiv:2509.21636v1 of this article is distributed under the Creative Commons Attribution 4.0 International licence, as recorded in the arXiv licence metadata for this exact version and shown on the abstract page https://arxiv.org/abs/2509.21636v1. A full rights notice is delivered at licenses/arxiv-2509.21636-CC-BY-4.0.txt. Characters: every retained excerpt is pure ASCII, so the delivered engine, XeLaTeX with its default Latin Modern text font, reports no missing character.
\begin{lem}\label{lem: embedding bimodules}
Let $U:\R\to \cO(\oplus_{i=1}^3 \sH_\R^{(i)})$ be an almost periodic representation with each $\sH_\R^{(i)}$ nontrivial.
Set $M=\Gamma_q(\oplus_{i=1}^3\sH_\R^{(i)}, U)''$ and $\chi$ its $q$-quasi free state.
For each subset $I\subset \{1,2,3\}$, let $M_I=\Gamma_q(\oplus_{i\in I} \sH_\R^{(i)}, U)''\subset M$
	and put $\cM=M\rtimes_{\sigma^\chi}\R$ with semifinite trace $\Tr_\chi$,
		and $\cM_I=M_I\rtimes_{\sigma^\chi}\R\subset \cM$.

%Let $U_i: \R\to \cO(\sH_\R^{(i)})$ be almost periodic orthogonal representations for $i=1,2,3$.
%For each subset $I\subset \{1,2,3\}$, set $M_I=\Gamma_q(\oplus_{i\in I} (\sH_\R^{(i)}, U_i))''$, $\chi_I$ the corresponding $q$-quasi free state,
%	and $\cM_I=\Gamma_q(\oplus_{i\in I} H_\R^{(i)}, \oplus_{i\in I} U_i)''\rtimes_{\sigma^{\chi_I}}\R$
%	with the corresponding semi-finite trace $\Tr_I$. 
%Pick $p\in L\R$ a nonzero finite trace projection and put $N=p\cM p$ and $N_I=p \cM_I p$.

Let $p\in L_\chi\R$ be a nonzero finite trace projection and set $N=p\cM p$, $N_I=p\cM_I p$.
Then we may identify $L^2(N_{\hat 3 })\ot_{ N_{1} } L^2(N_{\hat 2})$ as an $N_1$-$N_1$ sub-bimodule of $L^2(N)$,
	where we denote by $\hat i$ the set $I\setminus \{i\}$ and by $i$ the set $\{i\}$.
\end{lem}

\begin{proof}
Define $\theta: N_{\hat 3}\ot_{\rm alg} L^2(N_{\hat 2})\ni a\ot \xi \mapsto a\xi\in L^2(N)$,
	where we view $L^2(N_{\hat 2})$ as a subspace of $L^2(N)$.
It is clear that $\theta$ is $N_1$-bimodular 
	and thus it suffices to show that $\theta$ extends to an embedding on $L^2(N_{\hat 3 })\ot_{ N_{1} } L^2(N_{\hat 2})$.

	
For $\xi, \eta\in L^2(N_{\hat 2})$ and $a,b\in N_{\hat 3}$, one has
	$\langle \theta(a\ot \xi), \theta(b\ot \eta)\rangle=\langle a\xi, b\eta\rangle=\langle b^*a \xi, \eta\rangle$.
Thus to see $\theta$ extends to an isometry on $L^2(N_{\hat 3 })\ot_{ N_{1} } L^2(N_{\hat 2})$, 
	it is enough to prove that 
	$\langle x \xi, \eta\rangle=\langle E_{N_{1}}(x)\xi, \eta\rangle$ for any $x\in N_{\hat 3}$ and $\xi,\eta\in L^2(N_{\hat 2})$.

We claim that for $x\in C_c(\R, M_{\hat 3})$ and $\xi,\eta\in C_c(\R, \cF_q(\sH^{(1)}\oplus \sH^{(3)}))$, 
	one has
$$
\langle x\xi, \eta\rangle=\langle E_{\cM_1}(x)\xi,\eta\rangle,
$$
where $E_{\cM_1}: \cM\to \cM_1$ is the normal expectation.
Observe that from our claim one also has the same equation holds for $x\in p\cM_{\hat 3} p$
	and $\xi,\eta \in L^2(p\cM_{\hat 2} p)$ by a density argument.
Since ${E_{\cM_1}}_{\mid p\cM_1 p}$ coincides with the normal expectation from $p\cM p$ to $p\cM_1 p$,
	our desired conclusion follows.

To see our claim, we compute
\[
\begin{aligned}
\langle x\xi, \eta\rangle=\int_\R \langle \lambda_sx(s)\xi,\eta\rangle ds
	=\int_\R \Big(\int_\R \langle \sigma_{-t}(x(s))\xi(t),\eta(s+t)\rangle dt \Big)ds.
\end{aligned}
\]
If one has $\langle \sigma_{-t}(x(s))\xi(t),\eta(s+t)\rangle=\langle E_{M_1}\big (\sigma_{-t}(x(s))\big)\xi(t), \eta(s+t)\rangle$,
	then it follows that 
$$
\int_\R \langle \lambda_sx(s)\xi,\eta\rangle ds=\int_\R \langle \lambda_s E_{M_1}(x(s))\xi,\eta\rangle ds=\langle E_{\cM_1}(x)\xi,\eta\rangle.
$$
Therefore, we only need to prove that for $x\in \Gamma_q(\sH_\R^{(1)}\oplus \sH_\R^{(2)}, U)''$, $\xi,\eta\in \cF_q(\sH^{(1)}\oplus \sH^{(3)})$,
	one has $\langle x\xi,\eta\rangle=\langle E_{M_1}(x)\xi,\eta\rangle$, for which one easily verifies by using Wick formulas.
\end{proof}

\end{document}
